\begin{solapartat}
$N(t)$ és el nombre de nuclis que queden de la substància radioactiva, després que hagi passat un temps $t$, si en tenim un nombre $N_0$ a l'instant inicial. \punts{0,2}

L'exponent $0,005\ \mathrm{s^{-1}}$ és la constant de desintegració radioactiva $\Rightarrow \lambda = 0,005\ \mathrm{s^{-1}}$ \quad\punts{0,2}

El període de semidesintegració és el temps que ha de passar perquè es redueixi a la meitat la quantitat d'una substància radioactiva \punts{0,2} $\Rightarrow$
\begin{align*}
&\frac{N_0}{2} = N_0\,\mathrm{e}^{-\lambda\, t_{1/2}} \Rightarrow \frac{1}{2} = \mathrm{e}^{-\lambda\, t_{1/2}} \Rightarrow \ln(1) - \ln(2) = -\lambda\, t_{1/2} \Rightarrow\\
&t_{1/2} = \frac{\ln 2}{\lambda} = 1,4\cdot 10^{2}\ \mathrm{s} \quad\punts{0,4}
\end{align*}
\end{solapartat}

\begin{solapartat}
L'activitat d'una substància radioactiva es defineix com:
\[ A(t) = -\frac{\mathrm{d}N(t)}{\mathrm{d}t} = (-\lambda)(-N_0)\,\mathrm{e}^{-\lambda t} \quad\punts{0,5} \Rightarrow \]
\begin{align*}
A(t = 4\ \text{hores}) &= 0,005\ \mathrm{s^{-1}}\cdot 10^{28}\ \text{nuclis}\ \mathrm{e}^{-\left(0,005\ \mathrm{s^{-1}}\cdot 4\text{hores}\cdot 3600\frac{\mathrm{s}}{1\ \text{hora}}\right)}\\
&= 2,7\cdot 10^{-6}\,\frac{\text{nuclis}}{\mathrm{s}} \quad\punts{0,5}
\end{align*}
\end{solapartat}
